\chapter{Verbal Morphology and Inflection}\label{ch:verbs}\label{ch:verbal-morphology}

\section{The finite verbal system}

The verb is the morphological nucleus of the Iridian clause, and upon it falls
the principal burden of grammatical expression. Where nouns inflect for case and number, verbs
express the event's aspectual contour, the configuration of its participants,
and the speaker's modal orientation. Every finite verb bears the aspect-alignment suffix,
which is the sole obligatory element of the verbal complex; beyond this the verb
may carry markers of voice, mood, and evidentiality, in an order fixed by the
morphological template set out in the following section. Tense, gender, number
agreement, and pronominal cross-reference find no expression on the Iridian
verb.

In the schematic examples, double brackets enclose lexical meanings whose
Iridian stems are unspecified, while subscripts identify complete inflectional
cells (§\,\ref{sec:lexical-slot-schemata}). Thus
\ird{[[write]]}\textsubscript{\Prog{}} denotes a complete active progressive,
including its class-selected stem grade and aspect exponent; it does not
specify how an unlisted root forms that cell. Case and orientation remain
explicit wherever the contrast depends on them.

\section{The morphological template}\label{sec:morph-template}

The structure of a finite verb in Iridian is given by the following template,
wherein parentheses enclose optional elements and the sole obligatory position
is that of the aspect-alignment suffix:

\medskip
\begin{center}
(\Neg{}) + \Stem{} + (\Antip{} / \Appl{}) + (\Caus{}) + (\textsc{orientation}) + \synhead{Asp.Align} + (\textsc{mood}) + (\textsc{evid})
\end{center}
\medskip

The \synhead{Asp.Align} slot cumulatively realises the aspect and alignment
features of a single head and is the one element no finite verb may lack.
The mood slot may be analysed as containing a null morpheme in the indicative. The negative prefix \ird{za-}
stands before the entire stem complex. The indirect evidential has default
\ird{-al} after a consonant and contextual \ird{-l} after a vowel;
it follows finite aspect and mood and precedes the linker
(§\,\ref{sec:vm-evid}). Person is not morphologically expressed on the Iridian
verb: there is no pronominal clitic series and no agreement layer. Overt person
reference is conveyed by full pronoun DPs, demonstratives in the third person,
and zero anaphora where discourse recoverability licenses argument omission.

The \Stem{} position admits a derived lexical stem as well as a simple root.
The productive stative derivation and its inchoative extension belong inside
this position: ROOT--STATE--BECOME precedes inner valency, Caus and orientation
(§\,\ref{sec:stative-inchoative}). They do not reorder those outer heads.
The lexical preverbs \ird{de-}, \ird{mně-}, \ird{pře-},
\ird{le-}, \ird{ra-}, \ird{vra-}, \ird{ša-}, \ird{na-},
\ird{ru-}, \ird{po-}, and marginal \ird{ne-} belong inside that stem, as in
\ird{mněkař-} \trsl{read through} and \ird{lenočk-} \trsl{run out},
before the valency and orientation suffixes
(§\,\ref{sec:lexical-preverbs}). They modify the lexical event description
without supplying an Asp.Align value or changing the alignment domain;
indicative \ird{za-} attaches outside the complete derived stem.


\subsection{The Mirror Principle and functional projections}
\label{sec:vm-overview}
\label{sec:vm-mirror}

\noindent Iridian's causative suffix precedes outer patient orientation,
which in turn precedes finite aspect, as in \ird{pjášt-r-an-ik}
\trsl{cause to eat}, with perfective viewpoint and patient selection.
The corresponding analysis places Caus below orienting Voice and Asp.Align;
their distinct contributions to argument structure, participant selection and
viewpoint motivate this hierarchy alongside the suffix order
(§§\,\ref{sec:vm-selection-domains}, \ref{sec:vm-caus}).

Baker formulates the \term{Mirror Principle} as the requirement that
``Morphological derivations must directly reflect syntactic derivations
(and vice versa)'' \cite[p.~375]{baker1985}. The correspondence concerns ordered
operations and their grammatical consequences, rather than an identity between
an affix's position in a string and a head's position in a tree. Baker's
argument compares morphological and syntactic derivations with independently
specified content; his discussion of clitics and nonconcatenative morphology
also limits a direct inference from visible segments to derivational structure
\cite[§§2, 5]{baker1985}.

The Iridian analysis adopts a more limited working expectation: where
successive suffixal attachment realises successive syntactic combination,
inner morphological constituents correspond to lower syntactic domains.
Surface order constrains this analysis without uniquely determining it,
since head movement, affix-specific directionality and morphological merger
can mediate the correspondence \cite{harley2011,harley2013b}.
Table~\ref{tab:vm-projections} therefore distinguishes syntactic height from
the placement of an exponent within the verbal word. Verbal categorisation,
actor introduction and outer orientation occupy distinct positions; the
separation of verbalising $v$ and argument-introducing Voice has independent
motivation in Harley's analysis, while their particular ordering and licensing
conditions here belong to Iridian \cite{harley2013a,harley2017}.

\begin{table}
  \footnotesize
  \caption{Verbal projections in ascending syntactic order and their
  morphological realisation. Optional layers occur only where licensed.}
  \label{tab:vm-projections}
  \begin{tabularx}{\textwidth}{L l L l}
    \toprule
    {\sc syntactic projection} & {\sc exponent} & {\sc realisation} &
    {\sc see} \\
    \midrule
    $\surd$P / \synhead{VP} & \Stem{} & Lexical stem, including licensed preverbs
                 & §\,\ref{sec:vm-root} \\
    \synhead{vP} & $\varnothing$ & Categorisation of a simple lexical root
                 & §\,\ref{sec:vm-root} \\
    \synhead{VoiceP}$_{\mathrm{AP}}$ / \synhead{ApplP}
                 & \Antip{} / \Appl{} & Inner valency suffix; alternative heads
                 & §\,\ref{sec:vm-antipassive}, §\,\ref{sec:vm-appl} \\
    Inner actor layer & $\varnothing$ & Introduces the lexical actor where licensed
                 & §\,\ref{sec:vm-root} \\
    \synhead{CausP} & \Caus{} & Suffix outside inner valency
                 & §\,\ref{sec:vm-caus} \\
    \synhead{VoiceP}$_{\mathrm{orient}}$
                 & \ird{-n-}/\ird{-an-}, $\varnothing$ & Orientation outside Caus
                 & §\,\ref{sec:vm-n} \\
    \synhead{AspP} & \synhead{Asp.Align} & Aspect--alignment suffix outside orientation
                 & §\,\ref{sec:vm-asp} \\
    \synhead{NegP} & \Neg{} & Left-edge verbal proclitic in the indicative
                 & §\,\ref{sec:vm-neg} \\
    \synhead{MoodP} & \textsc{mood} & Suffix outside Asp.Align, or zero
                 & §\,\ref{sec:vm-moodp} \\
    \synhead{EvidP}$_\text{inner}$ & \ird{-al}/\ird{-l} & Suffix outside finite aspect and mood
                 & §\,\ref{sec:vm-evid} \\
    \bottomrule
  \end{tabularx}
\end{table}

\noindent Zero exponence leaves a syntactic position unpronounced;
cumulative exponence realises several features together; conditioned allomorphy
selects a form within an established morphological environment. None of these
operations by itself reverses morphological constituency. Prefixal \Neg{}
has a separate left-edge realisation instruction (§\,\ref{sec:vm-neg}), while
lexical preverbs belong inside the stem (§\,\ref{sec:lexical-preverbs}).
Asp.Align is one feature-bearing head throughout this analysis, without a
separate Align projection.

\section{The verb stem and thematic classes}

Iridian verbs admit no citation form. Lexicographers record the bare stem with a
following hyphen to indicate its character as a bound form requiring at least
one suffix; the verb rendered in English as \textit{to eat} is thus entered as
\ird{pjašt-}, where \ird{t} constitutes the \term{thematic consonant}.

The thematic consonant, being the final consonant of the stem, is phonologically
active: it undergoes regular alternations when certain suffixes are adjoined.
Suffixes select distinct mutation classes: antipassive \ird{-aš-} triggers
\term{sibilant replacement}, whereas simultaneous \ird{-ěc} and conjunctive
\ird{-e} select the front grade; irrealis \ird{-ěv} selects the
mixed grade (§\,\ref{sec:suffix-classes}). The grade is a property of the
construction rather than an automatic effect of the following vowel.
The five thematic classes in Table~\ref{tab:verbclasses} govern sibilant
replacement only; a listed stem shape does not independently establish that
the lexeme permits the antipassive argument frame.

\begin{table}
	\footnotesize
	\caption{Thematic consonant classes.}\label{tab:verbclasses}
	\begin{tabularx}{\textwidth}{l L L L}
		\toprule
		{\sc class} & {\sc thematic consonants} & {\sc change before \Antip{}} & {\sc example} \\
		\midrule
		I    & -t, -k, -c, -č  & → -č                    & \ird{pjašt-} → \ird{pjášč-} \\
		II   & -s, -š          & → -š                    & \ird{tras-} → \ird{tráš-} \\
		II-A & -l, -p          & add -š                  & \ird{dal-} → \ird{dalš-} \\
		III  & -d, -h, -j, -z, -ž & → -ž               & \ird{vad-} → \ird{váž-} \\
		III-A & all remaining  & add -ž                  & \ird{ščen-} → \ird{ščenž-} \\
		\bottomrule
	\end{tabularx}
\end{table}

Loan stems with surface \orth{g} are assigned analogically to the voiced
velar/fricative series and do not require a native phoneme /g/. Palatalising
suffixes do \emph{not} follow Table~\ref{tab:verbclasses}; they invoke the
consonantal-grade system described in
Section~\ref{sec:suffix-classes}.

Verbal stem grades are selected by the morphological construction and the lexical
class of the stem. Fleeting-vowel stems have a reduced grade, as in
\ird{razeň-}/\ird{razn-}; some modal constructions select a raised grade,
as in raised-strong \ird{zděk-}/\ird{zďíč-}; the perfective suffix has its own
\ird{-ek}/\ird{-ik-} relation. These changes preserve quantity unless an
independent quantity grade is specified. Consonant mutation applies to the stem
or aspect suffix selected in that domain, followed by syllabification and the
ordinary voicing rules (§\,\ref{sec:morphophon}).

\subsection{Quantitative classes and local grades}\label{sec:verbal-quantity}

Verbal quantitative class determines where an aspectual contrast receives
its quantity. The grades of §\,\ref{sec:morphophon} distinguish lexical
G0, strong G1, raised G2, and raised-strong G3. A grade selects one
licensed nucleus in the local stem, rather than lengthening every vowel
in a written word. The target ordinarily is the rightmost eligible
root nucleus, excluding fleeting vowels and material belonging to a
prefix or suffix; a lexical target specification takes precedence.
In \ird{zemac-} \trsl{run}, the target is its second vowel, /a/:
G1 gives \ird{zemác-}, preserving the initial /ɛ/ and final /ts/.
The digraph \orth{pj} in \ird{pjašt-} represents /pʲ/ and contains no
vowel nucleus; its strong grade is \ird{pjášt-}.

Class~A expresses the strong progressive and perfective in its root
grade, while the continuous, contemplative, and retrospective select
G0. Class~B instead uses strong aspectual exponents with an otherwise
short root. Class~C preserves an independently long lexical root
and selects short aspectual endings. The principal parts in
Table~\ref{tab:verbal-quantity-classes} illustrate these specifications;
the perfectives of \ird{zemac-} and \ird{knář-} are intransitive,
whereas ordinary transitive \ird{pjašt-} includes patient orientation
\ird{-n-}.

\begin{table}[htbp]
\footnotesize
\caption{Principal parts of the three verbal quantitative classes.}
\ifparadigmappendix\else\label{tab:verbal-quantity-classes}\fi
\begin{tabularx}{\textwidth}{L L L L L L}
\toprule
{\sc class/stem} & {\sc continuous} & {\sc progressive} & {\sc contemplative} & {\sc perfective} & {\sc retrospective} \\
\midrule
A \ird{zemac-} & \ird{zemaca} & \ird{zemácime} & \ird{zemacáž} & \ird{zemácek} & \ird{zemacáv} \\
B \ird{pjašt-} & \ird{pjašta} & \ird{pjaštíme} & \ird{pjaštáž} & \ird{pjaštník} & \ird{pjaštnáv} \\
C \ird{knář-} & \ird{knářa} & \ird{knářime} & \ird{knářaž} & \ird{knářek} & \ird{knářav} \\
\bottomrule
\end{tabularx}
\end{table}


Class~C \ird{strád-} likewise retains its lexical long vowel in
\ird{stráda}, \ird{strádime}, and \ird{strádaž}.

Copular, existential, and other functional predicates retain their listed
principal parts; surface shortness alone does not assign them to Class~A
(Chapter~\ref{ch:predicates}).

The Class~B perfective has long \ird{-ék}, with bound \ird{-ík-},
where a bare perfective is independently licensed, and long \ird{-ník}
with ordinary patient orientation. Its progressive has \ird{-íme}
and its retrospective \ird{-áv}, or \ird{-náv} after orientation.
Class~A instead combines root G1 with short perfective
\ird{-ek}/\ird{-nik} and short progressive \ird{-ime}; its
retrospective has short-root G0 and long \ird{-áv}/\ird{-náv}.
The stative verbaliser's perfective \ird{-ik} likewise follows the
stem's quantitative class: \ird{múžnik} has root G1 and short /ɪ/.
In the principal variety, free \ird{-ék} and bound \ird{-ík-} both
have /iː/; the traditional height relation remains audible in conservative
regional speech. Their morphological distribution and the linker-selected
consonantal grade remain distinct (§\,\ref{sec:vowel-raising}).

Valency derivation establishes a new local quantity domain before
Asp.Align attaches. Replacement antipassive Voice selects a strong
stem such as \ird{pjášč-}; compact causative \ird{-r-} selects G1
on an eligible local stem, as in \ird{pjáštr-}; and applicative
\ird{-éb-} retains its own long vowel. Such stems select short
adjacent aspectual exponents: \ird{pjášček}, \ird{pjáščav},
\ird{pjáštrime}, \ird{pjaštébaž}. A weak aspect does not undo
strength already supplied by the inner head. Additive antipassives,
including \ird{dalš-}, \ird{kanž-}, and \ird{lovž-}, preserve their
listed short root and short \ird{-ek}/\ird{-av}; Aspect does not
reach through that Voice head to create a root-strong grade.

The short \ird{-až} is a conditioned contemplative allomorph after
Class~C or a locally quantity-bearing valency stem. Short
\ird{-av}, \ird{-ime}, and \ird{-ek}/\ird{-ik} likewise belong
to their licensed local combinations. These choices do not impose
one long vowel per morphological word. Directive Mood, participant
nominalisers, relatives, and outer anchoring preserve their independently
specified quantities unless their construction explicitly selects
another allomorph. Thus \ird{pjáštrím} may contain both the
causative-selected root quantity and the long imperative vowel.

Morphological height selection remains independent of quantity. Raising under G2
preserves the selected vowel's quantity, while an affirmative irrealis
with an exposed, eligible root can select raised-strong G3. Mid
\ird{e/ě} then gives \ird{í}, and \ird{o} gives \ird{ú}, where
that raised relation is licensed; low-vowel \ird{pjašt-} selects
strong \ird{pjášč-} without an invented raising of /a/.
For eligible mid nuclei, G1 and G3 have the same long-vowel outputs
/iː, uː/ in the principal variety, although their selection and traditional
spellings remain distinct. The consonantal grade is interpreted separately
at its own boundary.
Ordinary indicative Aspect may preselect its class-assigned root grade
before assembly with plain patient orientation \ird{-n-}, as in
Class~A \ird{vádnik}. This construction-specific selection does not
make orientation transparent to every later head. Mood instead acts on
the assembled local domain: intervening valency or orientation prevents
it from reaching back to assign a new root grade
(§\,\ref{sec:vm-moodp}).


\section{Aspect and alignment}
\label{sec:aspect}

Iridian grammaticalises \term{aspect} rather than \term{tense} on the finite
verb. There is no dedicated past-, present-, or future-tense suffix; temporal
location is inferred from discourse context, temporal adverbials, and the
lexical semantics of the predicate. Aspect and alignment are fused into a single
obligatory suffix, the \term{aspect-alignment morpheme} (Asp.Align), which every
finite verb must bear: selecting an aspectual category simultaneously determines
how the event is construed and the alignment domain within which
valency and orientation determine the privileged participant.

The five \synhead{Asp.Align} categories fall into two domains. The
\term{perfective domain} comprises the perfective and retrospective, which
present an eventuality as bounded or anterior and license the ordinary
ergative-absolutive frame: patient direct, agent transitive. The
\term{imperfective domain} comprises the continuous, progressive, and
contemplative, which present an eventuality internally or prospectively and
license the ordinary nominative-accusative frame: agent direct, patient
indirect. Applicatives, antipassives, and patient-selecting Voice modify these
defaults (§\,\ref{sec:vm-asp-case}). The case split belongs to Iridian's
grammar; temporal viewpoint alone does not determine the selected participant
or confer aboutness topichood.

A transitive writing predicate makes the case alternation visible while
holding the participants and lexical event constant:

\pex\label{ex:vm-slot-alignment}
\a
\begingl
\gla Veň \textnormal{[[letter]]}\textsubscript{\Ind{}} \textnormal{[[write]]}\textsubscript{\Prog{}}.//
\glb woman.\Dir{} letter.\Ind{} write-\Prog{}//
\glft \trsl{A woman is writing a letter.}
\lingcontext{ordinary active; writing in progress}//
\endgl
\a
\begingl
\gla \textnormal{[[letter]]}\textsubscript{\Dir{}} veňa \textnormal{[[write]]}\textsubscript{\Trs{}.\Pf{}}.//
\glb letter.\Dir{} woman.\Trs{} write-\Trs{}-\Pf{}//
\glft \trsl{A woman wrote a letter.}
\lingcontext{ordinary patient orientation; one completed letter}//
\endgl
\xe

The woman remains the actor in both clauses, although only the progressive
assigns that participant direct case. In the perfective, the letter takes direct
case and the woman transitive case; \Trs{}.\Pf{} includes ordinary
patient orientation and the perfective exponent, without assigning the
unlisted stem a particular quantitative class.

The account below first lays out the shared temporal-semantic machinery
(including Reichenbachian framing and the formal Asp.Align analysis), then treats
the perfective domain (§\,\ref{sec:aspect-pf}) before turning to the imperfective domain
in §\,\ref{sec:aspect-ipf}.

\subsection{Viewpoint and temporal interpretation}

Three temporal notions must be kept distinct. Tense locates a situation in time
relative to a reference point, most commonly the moment of speaking. Grammatical
aspect does not locate the situation in this way, but construes its internal
temporal contour: whether it is viewed as bounded, ongoing, prior from a later
vantage point, or not yet realised. \term{Lexical aspect} or \term{Aktionsart},
by contrast, belongs to the predicate itself. A predicate may denote a state,
activity, accomplishment, or achievement independently of the inflection it
bears. Iridian aspectual morphology therefore does not determine the inherent
temporal type of the eventuality; rather, it imposes a viewpoint upon an
eventuality whose lexical contour is supplied by the verb phrase.

Lexical preverbs make this distinction overt: \ird{mněkař-} \trsl{read
through} specifies traversal of a text, while its continuous and progressive
forms describe that traversal without asserting arrival at its endpoint.
The same stem admits all five Asp.Align categories, with the ordinary case
frames and Voice requirements (§\,\ref{sec:preverb-aspect}). Likewise,
\ird{ranóčkime} \trsl{is running into an enclosure} describes an ingress
in progress without asserting that entry has been completed, while
\ird{rustójime} \trsl{is retracing a route} locates the returning motion
within its unfolding course. Direction, configuration, and restitution belong
to the lexical description; continuous, progressive, contemplative, perfective,
and retrospective morphology remain available wherever the resulting meaning
licenses them.

The difference between an unfolding change and its viewpoint is visible in
a river-widening situation. Smith, Hopperdietzel and Koontz-Garboden's
Japanese example contains a lexical inchoative with present morphology;
the corresponding Iridian schema places progressive aspect outside the
derived change predicate:

\pex\label{ex:vm-comparative-widening}
\a
\begingl
\glpreamble \textit{Japanese} \cite[example~9b, p.~332]{smithhopperdietzelkoontzgarboden2025}//
\gla kawa-ga osoku hiro-gar-u.//
\glb river-\textsc{nom} slowly be.wide-\textsc{inch}-\textsc{prs}//
\glft \trsl{The river is slowly widening.}//
\endgl
\a
\begingl
\glpreamble \textit{Iridian; schematic}//
\gla \textnormal{[[river]]}\textsubscript{\Dir{}} \textnormal{[[slowly]]}\textsubscript{\textsc{adv}} \textnormal{[[wide]]}\textsubscript{\textsc{state.inch}.\Prog{}}.//
\glb river.\Dir{} slowly.\textsc{adv} wide-\textsc{state}-\textsc{inch}-\Prog{}//
\glft \trsl{The river is slowly widening.}
\lingcontext{erosion is increasing its width at the reference point; no final width is asserted}//
\endgl
\xe

In the Iridian predicate, the inchoative supplies the change whose rate the
manner expression modifies, while the progressive gives an internal view of
that change. Replacing a state predicate's continuous with its progressive
does not itself supply the inchoative layer
(§\,\ref{sec:stative-inchoative}). The shared English translation of
(\ref{ex:vm-comparative-widening}a--b) thus conceals a difference between
Japanese present morphology and Iridian viewpoint morphology.

A useful descriptive frame is Reichenbach's distinction among speech time
(\textit{S}), reference time (\textit{R}), and event time (\textit{E}):
\textit{S} is the moment of utterance, \textit{E} the time at which the
situation itself takes place, and \textit{R} the temporal vantage point from
which that situation is presented. The usefulness of \textit{R} is that speakers
do not always describe an event directly from the moment of speaking, but may
instead locate it from some later or earlier point already established in the
discourse; English \textit{had left}, for example, places the leaving
(\textit{E}) before a reference point (\textit{R}) that is itself prior to the
speech moment (\textit{S}). In a tense-prominent system, finite morphology often
fixes the relation among all three. In Iridian, by contrast, the finite suffix
primarily construes the relation between \textit{E} and \textit{R}, while the
relation of \textit{R} to \textit{S} is often left to context, temporal
adverbials, and discourse sequencing. This is why the same aspectual category
may receive past-, present-, or future-oriented readings in different
environments without ceasing to be aspectual rather than temporal.

The same progressive can describe a ferry approaching now or at a
past reference point supplied by a narrative:

\pex\label{ex:vm-slot-reference-time}
\a
\begingl
\gla \textnormal{[[ferry]]}\textsubscript{\Dir{}} \textnormal{[[approach]]}\textsubscript{\Prog{}}.//
\glb ferry.\Dir{} approach-\Prog{}//
\glft \trsl{A ferry is approaching.}
\lingcontext{a lookout reports movement through the harbour fog at speech time; $R=S$}//
\endgl
\a
\begingl
\gla \textnormal{[[ferry]]}\textsubscript{\Dir{}} \textnormal{[[approach]]}\textsubscript{\Prog{}}.//
\glb ferry.\Dir{} approach-\Prog{}//
\glft \trsl{A ferry was approaching.}
\lingcontext{a narrative has established the previous afternoon, before the harbour closed; $R<S$}//
\endgl
\xe

Both clauses locate the reference point within the approaching event,
without asserting that the ferry reaches the harbour. The English past in
(b) renders the narrative location of that point, while the Iridian
predicate retains the same progressive cell as in (a).

Japanese and Korean also distinguish the temporal anchor from the
aspectual interpretation: Sato and Cho's comparison separates absolute
main-clause tense from subordinate temporal reference and examines
existential-auxiliary aspectual constructions \cite{satocho2024aspect}.
Iridian's unchanged progressive in (\ref{ex:vm-slot-reference-time})
illustrates a different grammatical division of labour, since discourse
locates the reference point without a morphological tense opposition.

The five categories may be sketched in these terms. The \term{perfective}
presents an eventuality as bounded and viewed as a whole, with \textit{R}
aligned to the event in its entirety. The \term{retrospective} places \textit{E}
before \textit{R}, yielding a later vantage point from which present relevance,
result, prior experience, or other discourse effects of anteriority may be
construed. The \term{continuous} and \term{progressive} both place \textit{R}
within \textit{E}, but they differ in the kind of internal view imposed: the
continuous presents the situation as simply obtaining at reference time, without
foregrounding change, onset, or unfolding contour, whereas the progressive more
narrowly profiles a dynamic event as actively unfolding at the reference point.
The \term{contemplative} places \textit{E} after \textit{R}, presenting the
event as projected, intended, expected, or impending rather than as obtaining or
complete.

These five categories divide naturally into two families. The
\term{perfective domain}---the perfective and retrospective---consists of those
categories that locate the event from an external or posterior vantage and
present it as bounded or anteriorly packaged: in Reichenbachian terms, the
event time is enclosed within the reference span, or precedes it. The
\term{imperfective domain}---the continuous, progressive, and
contemplative---consists of those categories that do not present the event as a
bounded whole, but instead either locate reference time within the ongoing event
or situate reference time prior to an unrealized projection of the event. The
Reichenbachian figure (Figure~\ref{fig:iridian-aspect-reichenbach}) makes this
grouping visually explicit: the upper two rows belong to the perfective domain,
in which the event is viewed as bounded or anterior from outside; the lower
three rows belong to the imperfective domain, in which R lies within E or
precedes an unrealized E.

Because viewpoint aspect interacts with lexical aspect, the same inflection does
not yield identical interpretive effects with every predicate class. With states
and activities, the continuous naturally yields persistent or ongoing readings;
with accomplishments it typically selects an internal phase rather than a
completed endpoint; with achievements it may invite iterative, preparatory, or
inceptive coercions. Conversely, the perfective and retrospective readily
package telic predicates as bounded wholes or prior results, but with atelic
predicates they instead delimit a stretch, occasion, or state viewed from an
external vantage. The exact semantic and distributional consequences of these
interactions are discussed in the domain subsections below.

Typologically, this pattern is an instance of aspect-conditioned split
ergativity. Iridian does not maintain a single alignment pattern across all
finite clauses; rather, the aspectual category selected on the verb determines
whether the clause is built on an ergative-absolutive or a nominative-accusative
frame. Aspect in Iridian therefore has consequences at two levels at once: it
construes the temporal profile of the eventuality and it determines the
morphosyntactic alignment through which the clause is organised.

The aspectual domains condition case without encoding a discourse topic.
A perfective transitive may continue reference to its agent, although the
patient takes direct case, while an imperfective clause may introduce an event
without an established topic. Valency and orientation morphology determine
which participant receives direct-case privileges within the aspectual domain;
linked topicalisation independently determines aboutness
(§\,\ref{sec:vm-asp-case}; Chapter~\ref{ch:infstr}).

\begin{figure}[t]
\centering
\begin{tikzpicture}[
  x=1.05cm,y=1cm,
  evt/.style={line width=1.1pt},
  refpt/.style={circle,fill,inner sep=1.3pt},
  speech/.style={diamond,draw,inner sep=1.2pt},
  rowlab/.style={anchor=east,font=\small},
  every node/.style={font=\small}
]

% All rows hold S at x = 4.5 so the deictic baseline is visually constant.

% perfective (y = 0)
\node[rowlab] at (0,0) {perfective};
\draw[->] (0.5,0) -- (8.5,0);
\draw[evt] (3.0,0) -- (6.0,0);
\draw (3.0,0.15) -- (3.0,-0.15);
\draw (6.0,0.15) -- (6.0,-0.15);
\node[refpt] at (3.0,0) {};
\node[refpt] at (6.0,0) {};
% R bracket above the timeline, E label between bracket and line
\node at (4.5,0.30) {$E$};
\draw[decorate,decoration={brace,amplitude=3pt}]
  (3.0,0.55) -- (6.0,0.55)
  node[midway,above=4pt] {$R$};
% S below, matching spacing of other rows
\draw[dashed] (4.5,-0.08) -- (4.5,-0.42);
\node[speech] (s1) at (4.5,-0.55) {};
\node[below=2pt of s1] {$S$};

% retrospective (y = -2.0)
\node[rowlab] at (0,-2.0) {retrospective};
\draw[->] (0.5,-2.0) -- (8.5,-2.0);
\draw[evt] (1.8,-2.0) -- (3.4,-2.0);
\draw (1.8,-1.85) -- (1.8,-2.15);
\draw (3.4,-1.85) -- (3.4,-2.15);
\node at (2.6,-1.65) {$E$};
\node[refpt,label=above:{$R$}] at (4.5,-2.0) {};
\node[speech] (s2) at (4.5,-2.55) {};
\node[below=2pt of s2] {$S$};
\draw[dashed] (4.5,-2.08) -- (4.5,-2.42);

% continuous (y = -4.0)
\node[rowlab] at (0,-4.0) {continuous};
\draw[->] (0.5,-4.0) -- (8.5,-4.0);
\draw[dotted] (1.6,-4.0) -- (2.2,-4.0);
\draw[evt] (2.2,-4.0) -- (6.8,-4.0);
\draw[dotted] (6.8,-4.0) -- (7.4,-4.0);
\node at (3.0,-3.65) {$E$};
\node[refpt,label=above:{$R$}] at (4.5,-4.0) {};
\node[speech] (s3) at (4.5,-4.55) {};
\node[below=2pt of s3] {$S$};
\draw[dashed] (4.5,-4.08) -- (4.5,-4.42);

% progressive (y = -6.2)
\node[rowlab] at (0,-6.2) {progressive};
\draw[->] (0.5,-6.2) -- (8.5,-6.2);
\draw[evt] (2.4,-6.2) -- (7.0,-6.2);
\draw (2.4,-6.05) -- (2.4,-6.35);
\draw (7.0,-6.05) -- (7.0,-6.35);
\node at (6.5,-5.85) {$E$};
\draw[->] (2.8,-5.55) -- (4.2,-5.55)
  node[midway,above,font=\scriptsize] {dynamic phase};
\node[refpt,label=above:{$R$}] at (4.5,-6.2) {};
\node[speech] (s4) at (4.5,-6.75) {};
\node[below=2pt of s4] {$S$};
\draw[dashed] (4.5,-6.28) -- (4.5,-6.62);

% contemplative (y = -8.2)
\node[rowlab] at (0,-8.2) {contemplative};
\draw[->] (0.5,-8.2) -- (8.5,-8.2);
\node[refpt,label=above:{$R$}] at (4.5,-8.2) {};
\draw[evt] (5.8,-8.2) -- (7.5,-8.2);
\draw (5.8,-8.05) -- (5.8,-8.35);
\draw (7.5,-8.05) -- (7.5,-8.35);
\node at (6.65,-7.85) {$E$};
\node[speech] (s5) at (4.5,-8.75) {};
\node[below=2pt of s5] {$S$};
\draw[dashed] (4.5,-8.28) -- (4.5,-8.62);

\end{tikzpicture}
\caption{Reichenbachian schematic for Iridian aspect. $E$ = event time, $R$ =
reference time, $S$ = speech time. In the perfective, $R$ spans the bounded
event (marked by dots at the endpoints and a bracket); elsewhere, $R$ is a
point-like vantage. The upper two rows (perfective, retrospective) belong to the
\term{perfective domain}: in both, the event is construed as bounded or
anterior, viewed from outside or from a posterior vantage. The lower three rows
(continuous, progressive, contemplative) belong to the \term{imperfective
domain}: in all three, $R$ either lies within $E$ or precedes an unrealized $E$,
and no external terminus is foregrounded. Dashed connectors mark the unmarked
deictic baseline in which $R$ coincides with $S$. $S$~is held at a constant
horizontal position across all rows; in discourse, adverbials or sequencing may
shift $R$ independently of $S$.}
\label{fig:iridian-aspect-reichenbach}
\end{figure}

\subsection{The Asp.Align system}\label{sec:vm-asp}

The obligatory aspect-alignment suffix encodes viewpoint and the alignment
domain within which the verbal construction licenses case. The five indicative
exponents are given in Table~\ref{tab:aspect}.

\begin{table}[ht]
	\footnotesize
	\caption{Aspect-alignment exponents by domain and quantitative class.}\ifparadigmappendix\else\label{tab:aspect}\fi
	\begin{tabularx}{0.98\textwidth}{L L L L}
		\toprule
		{\sc domain} & {\sc aspect} & {\sc exponent} & {\sc quantity and orientation} \\
		\midrule
		Perfective & Perfective & \ird{-ek}/\ird{-ék} & bare forms require an intransitive reading; ordinary transitive \ird{-nik}/\ird{-ník}; class-selected short or long bound \ird{-ik-}/\ird{-ík-} \\
		& Retrospective & \ird{-áv}/\ird{-av} & ordinary A/B long \ird{-áv}, or \ird{-náv} with orientation; C and listed strong-valency stems short \ird{-av}/\ird{-nav} \\
		\midrule
		Imperfective & Continuous & \ird{-a} & local G0 in ordinary A/B; previously selected lexical or valency strength retained \\
		& Progressive & \ird{-ime}/\ird{-íme} & A root G1 with short ending; B suffix-strong; C and strong-valency stem with short ending \\
		& Contemplative & \ird{-áž}/\ird{-až} & A/B ordinary G0 with long ending; C and strong-valency stem with short ending \\
		\bottomrule
	\end{tabularx}
\end{table}


Antipassive and applicative forms preserve the orientation and quantity
selected by their own constructions. Replacement antipassives such as
\ird{pjášč-} combine with short \ird{-ek}/\ird{-av}, while the
applicative \ird{-éb-} selects short perfective \ird{-ik}.
The distribution in Table~\ref{tab:verbal-quantity-classes} supplies the
quantity distinctions without changing the five viewpoint categories.

\subsection{Perfective and retrospective aspect}\label{sec:aspect-pf}

The perfective domain comprises those Asp.Align forms which construe an
eventuality as bounded and externally viewed. Unlike the imperfective domain,
which presents an event from within its internal course, the perfective domain
presents it as a temporally delimited whole. Iridian distinguishes two
categories within this domain: the \term{perfective} proper, which presents a
bounded eventuality as a discrete occurrence, and the \term{retrospective},
which presents a prior eventuality in relation to a present state, present
relevance, or present experience. Eventive bases supply events; simple
stative bases supply states under their independently licensed viewpoint
construals. Both categories interact regularly with
valency and alignment morphology, and both therefore belong to a single
structural and semantic domain rather than to unrelated aspectual classes.

Free perfective quantity follows the selected quantitative class.
Class~A \ird{kánek} \trsl{sang} and \ird{ščének} \trsl{arrived}
have root G1 and short \ird{-ek}; Class~B has long \ird{-ék} where
a bare perfective is licensed. Adjacent linker \ird{-e} and participant
nominaliser \ird{-ou} select the corresponding bound height grade,
short \ird{-ik-} or long \ird{-ík-}. A replacement antipassive
already supplies root strength: \ird{pjášček}, linked
\ird{pjáščice}, and participant \ird{pjáščikou} therefore
retain short suffix vowels (§\,\ref{sec:pf-dissim}).

Ordinary patient orientation combines \ird{-n-} with the class-selected
high-vowel perfective: short \ird{-nik} in root-strong Class~A and
long \ird{-ník} in suffix-strong Class~B. The stative verbaliser
\ird{-n-} of the root-strong property class independently selects short
\ird{-ik}, with root quantity assigned by that class; its surface
\ird{-nik} does not insert transitive Voice into the stative predicate.
Relational \ird{-(e)zn-} retains its separately specified quantitative
class: a suffix-strong short relational stem has long perfective
\ird{-ík}, whereas protected Class~C takes short \ird{-ik}
(Chapter~\ref{ch:derivational}). The selected quantities
persist before further morphology. Retrospective likewise distinguishes
ordinary A/B \ird{-áv}/\ird{-náv} from C and strong-valency
\ird{-av}/\ird{-nav}; antipassive \ird{pjáščav} has a strong
replacement stem and short ending.

The semantic opposition within the perfective domain is best stated as one of
\emph{event occurrence} versus \emph{post-event relevance}. The perfective
proper presents a bounded eventuality as a single occurrence situated at a definite
point in narrative or discourse time, and is therefore the primary narrative
aspect. With an eventive base, the retrospective presents a bounded prior event, but does
not profile that event merely as an occurrence; rather, it presents it through
its present relevance, present consequences, or experiential force. With a simple state, its anterior eventuality remains a state and no transition
is entailed. In this respect it corresponds in some uses to the English perfect. Its non-volitional
or accidental readings belong to the same general semantic profile: the event is
construed from the standpoint of its aftermath or present bearing, rather than
simply narrated as a bounded occurrence.

For this reason, speakers often choose the retrospective where reduced
volitionality is pragmatically inferable or interactionally relevant. By
presenting the event through its resulting situation rather than as a discrete
occurrence in the narrative line, the retrospective readily backgrounds agency
and favours readings such as `it happened', `it turned out that way', or `I find
myself in the state resulting from it'. This preference is pragmatic rather
than semantic: the retrospective does not itself encode non-volitionality, but
it is frequently selected in contexts where non-control, accident, or
unintended outcome is to be understood.

This contrast may be seen in the following pair:

\pex
\a
\begingl
\gla Praže de stožáv.//
\glb Prague.\Ind{} to go-\Ret{}//
\glft \trsl{(I) have been to Prague before.}//
\endgl
\a Compare:\par
\begingl
\gla Bych Praže de stóžek.//
\glb yesterday Prague.\Ind{} to go-\Pf{}//
\glft \trsl{Yesterday I went to Prague.} \lingcontext{specific past event; perfective}//
\endgl
\xe

Within the perfective domain, the ordinary transitive selects its patient
through \ird{-n-}, assigning direct case to that participant and transitive
case to the agent. An antipassive instead leaves only the actor independently
licensed, whereas a simple applicative selects its applied participant. An
inherently intransitive verb requires no overt orientation marker, since its
sole argument takes direct case (§\,\ref{sec:vm-asp-case}).

The retrospective can present a patient-related outcome in an ergative
transitive or an agent's prior experience in an antipassive. The latter has no
independent patient argument whose resulting state could supply the discourse
centre, although the activity's consequences may remain relevant. These are
contextual interpretations of the respective argument frames, as illustrated
by English \textit{the bread has been eaten} and \textit{I have eaten};
retrospective morphology itself does not require either a result-state or an
experiential reading.

With a transitive cooking predicate, the perfective can advance a narrative
by reporting the preparation of dinner, whereas the retrospective can present
the same event as relevant when someone asks whether dinner is ready:

\pex\label{ex:vm-slot-occurrence-anteriority}
\a
\begingl
\gla \textnormal{[[dinner]]}\textsubscript{\Dir{}} papky márnik.//
\glb dinner.\Dir{} father.\Trs{} cook-\Trs{}-\Pf{}//
\glft \trsl{A father cooked dinner.}
\lingcontext{the cooking is one completed event in an account of the evening}//
\endgl
\a
\begingl
\gla \textnormal{[[dinner]]}\textsubscript{\Dir{}} papky marnáv.//
\glb dinner.\Dir{} father.\Trs{} cook-\Trs{}-\Ret{}//
\glft \trsl{A father has cooked dinner.}
\lingcontext{someone asks whether dinner is ready; the prior cooking answers that question}//
\endgl
\xe

Both members retain patient-direct and agent-transitive case. The readiness
in (b) belongs to its context: retrospective morphology does not independently
guarantee that the food is still available or that anyone has eaten it.

Negation further confirms the semantic distinction between the two members of
the perfective domain. When negated, the perfective normally carries the
implicature that the event was expected, intended, or contextually projected,
but failed to occur. The retrospective, by contrast, is the neutral form for
simple past negation without this implication. The opposition is illustrated
below:

\pex
\a
\begingl
\gla Marek zapjášček.//
\glb Marek.\Dir{} \Neg{}-eat-\Antip{}-\Pf{}//
\glft \trsl{Marek didn't eat (but he should have).}//
\endgl
\a
\begingl
\gla Marek zapjáščav.//
\glb Marek.\Dir{} \Neg{}-eat-\Antip{}-\Ret{}//
\glft \trsl{Marek didn't eat (neutral past).}//
\endgl
\xe

An intransitive signal illustrates the contrast without introducing an
agent whose intentions might otherwise obscure the source of the expectation:

\pex\label{ex:vm-slot-negative-expectation}
\a
\begingl
\gla Spum zaraspúmitek.//
\glb bell.\Dir{} \Neg{}-ring.out-\Pf{}//
\glft \trsl{The bell did not ring out.}
\lingcontext{a scheduled signal failed to occur; the bell's ringing was expected}//
\endgl
\a
\begingl
\gla Spum zaraspumitáv.//
\glb bell.\Dir{} \Neg{}-ring.out-\Ret{}//
\glft \trsl{The bell has not rung out.}
\lingcontext{a neutral report of non-occurrence; no scheduled signal is presupposed}//
\endgl
\xe

The schedule supports the failed-expectation reading in (a), while the
negative retrospective in (b) reports anterior non-occurrence without
requiring that expectation. Both clauses negate the same lexical event,
\trsl{ring out}; neither form assigns volition to the bell.

\label{sec:ambitransitive}A small class of verbs shows particularly clearly how the perfective domain
interacts with argument structure. These are lexical ambitransitives: verbs
which possess both an intransitive reading, in which no patient is structurally
available, and a transitive reading, in which a patient is licensed. Neither
reading is derived from the other by ordinary valency operations; both are
lexical possibilities of the same stem. Two examples are \ird{lov-} \trsl{hunt,
forage [intr.]; catch, take [tr.]} and \ird{kan-} \trsl{sing, perform [intr.];
sing a song, chant a text [tr.]}.

The activity uses of \ird{pjašt-} \trsl{eat} and \ird{kař-}
\trsl{read} have a narrower distribution: a patientless lexical frame is
available in the indicative NOM domain, where it shares the finite form of
the transitive frame, but does not extend to an unmarked ERG-domain
intransitive. Thus \ird{Marek pjaštíme} may describe eating without a
referential patient, whereas the corresponding bounded activity requires
\ird{Marek pjášček}, with antipassive Voice. This aspect-conditioned
alternation differs from the freely available intransitive perfectives of
\ird{lov-} and \ird{kan-}; it also differs from leaving an independently
licensed patient unpronounced (§\,\ref{sec:nom-activity}).

For the lexical ambitransitives, the diagnostic for the two readings is
whether ordinary transitive orientation \ird{-n-} is available in one reading
but impossible in the other. In \ird{lov-}, the
intransitive perfective---formed without any Voice head, with the ERG Asp.Align
suffix attached directly to the stem---accesses only the intransitive reading:
the event is construed simply as hunting, with no patient structurally present
and no quarry included in the event frame as an affected participant. The
ergative and antipassive forms, by contrast, access the transitive reading. The
ergative perfective \ird{lóvnik} is possible only when something is caught or
taken; the antipassive construes the event as directed toward prey, but
suppresses that prey from the clause's core argument structure.

\pex
\a
\begingl
\gla Marek bych lóvek.//
\glb Marek.\Dir{} yesterday hunt-\Pf{}//
\glft \trsl{Marek went hunting yesterday.} \lingcontext{intransitive reading; no patient in event frame}//
\endgl
\a
\begingl
\gla Zvěrato Marku lóvnik.//
\glb game.\Dir{}.\Def{} Marek.\Trs{} hunt-\textsc{trs}-\Pf{}//
\glft \trsl{Marek caught the game.} \lingcontext{transitive reading; ergative}//
\endgl
\a
\begingl
\gla Marek bych lovžek.//
\glb Marek.\Dir{} yesterday hunt-\Antip{}-\Pf{}//
\glft \trsl{Marek was hunting yesterday.} \lingcontext{transitive reading; antipassive; quarry conceptually present but not licensed}//
\endgl
\xe

The two non-ergative forms are therefore not interchangeable. \ird{Lóvek}
belongs to the intransitive lexical reading: hunting is construed as an activity
with no patient present even in the event structure. \ird{Lovžek}, however,
belongs to the transitive lexical reading: the event is directed toward prey,
and that prey remains part of the conceptual frame even though the antipassive
removes it from core syntax. The distinction is thus not merely one of overt
argument expression, but of underlying event construal.

The stem \ird{kan-} likewise distinguishes an ordinary intransitive perfective,
a transitive perfective, and an antipassive:

\pex
\a \phon{kánek}{ˈkaː.nɛk}{sang, performed} \lingcontext{intransitive; no song in the event frame}
\a \phon{kánnik}{ˈkaːn.nɪk}{sang it, sang a specific song} \lingcontext{transitive; ergative; patient licensed}
\a \phon{kanžek}{ˈkan.ʒɛk}{was singing} \lingcontext{transitive reading; antipassive; patient suppressed from core syntax}
\xe

In \ird{kánek}, ordinary \ird{-ek} follows the full stem: G1 lengthens
/a/ to /aː/. Transitive \ird{kánnik} contains patient orientation
\ird{-n-} after the root-final \ird{n} and the construction-specific
high-vowel perfective. Both nasal consonants are retained across the
syllable boundary. The difference between \ird{e} and \ird{i} belongs to
perfective exponent selection. Antipassive \ird{kanžek} takes a short
free perfective after the Class~III-A augmented stem \ird{kanž-}.
Class~A root G1 supplies long \ird{á} in \ird{kánek} and \ird{kánnik},
whereas additive Voice in \ird{kanžek} retains short \ird{a}.
The forms distinguish valency and event construal together with their
licensed quantity patterns.

The transitive stem \ird{kař-} \trsl{read} also belongs to quantity
Class~A. Its lexical \ird{ř} represents /ʒ/; ordinary patient orientation
retains that consonant, as in \phon{kářnik}{ˈkaːʒ.nɪk}{read it}.
Antipassive replacement selects Class~III /ʒ/ with root G1, giving
\phon{kážek}{ˈkaː.ʒɛk}{was reading}. Its bound perfective is
\ird{kážik-}, as in participant-oriented \ird{kážikový}.
The replacement spelling \ird{ž} records the antipassive exponent;
it is independent of PAL (§\,\ref{sec:suffix-classes}).

\begin{table}[htbp]
\footnotesize
\caption{Class~A principal parts of \ird{kař-} and \ird{kan-}.}
\label{tab:read-sing-principal-parts}
\begin{tabularx}{\textwidth}{L L L L L L}
\toprule
{\sc stem} & {\sc continuous} & {\sc progressive} & {\sc contemplative} & {\sc perfective} & {\sc retrospective} \\
\midrule
\ird{kař-} & \ird{kařa} & \ird{kářime} & \ird{kařáž} & \ird{kářnik} & \ird{kařnáv} \\
\ird{kan-} (intr.) & \ird{kana} & \ird{kánime} & \ird{kanáž} & \ird{kánek} & \ird{kanáv} \\
\ird{kan-} (tr.) & \ird{kana} & \ird{kánime} & \ird{kanáž} & \ird{kánnik} & \ird{kannáv} \\
\bottomrule
\end{tabularx}
\end{table}

Direct linker suffixation selects bound perfective \ird{-ik-} and then
the front grade at its final /k/. The isolated and linked pairs are
\phon{kánek}{ˈkaː.nɛk}{sang} / \phon{kánice}{ˈkaː.nɪ.tsɛ}{sang, linked},
\phon{kánnik}{ˈkaːn.nɪk}{sang it} / \phon{kánnice}{ˈkaːn.nɪ.tsɛ}{sang it, linked},
and \phon{kanžek}{ˈkan.ʒɛk}{was singing} / \phon{kanžice}{ˈkan.ʒɪ.tsɛ}{was singing, linked}.
Likewise \ird{kař-} has
\phon{kářnice}{ˈkaːʒ.nɪ.tsɛ}{read it, linked} and
\phon{kážice}{ˈkaː.ʒɪ.tsɛ}{was reading, linked}.
The linker preserves the root and Voice grades
(§\,\ref{sec:conjunctive-linker}).

\section{Imperfective aspect}
\label{sec:aspect-ipf}

The imperfective domain comprises those Asp.Align categories that do not
construe the eventuality as a bounded whole viewed from outside. Where the
perfective domain packages an event as a closed unit and foregrounds its
culmination or its anteriorly packaged result, the imperfective domain takes
no such external vantage: it either locates the reference point within the
ongoing event, without presenting its boundaries as salient, or presents the
event as projected relative to reference time---not yet realised, but
anticipated, intended, or impending. These temporal properties characterise
the domain independently of its ordinary agent-direct case frame.

Three categories constitute the domain. The \term{continuous} is the unmarked
imperfective member; it presents a situation as simply obtaining at reference
time, without emphasis on change, onset, dynamic unfolding, or endpoint. The
\term{progressive} is the marked dynamic member; it profiles an event
specifically in its active unfolding phase. The \term{contemplative} is the
projective member; it presents the event as not yet realised at reference time
but projected from that point---intended, expected, impending, planned, or
otherwise prospectively construed. Together the three define an internal
opposition within the domain: static holding versus dynamic unfolding versus
prospective projection.

The three imperfective categories license the ordinary NOM-domain frame,
with an agent or sole participant in direct case and a retained patient in
indirect case. A simple applicative instead selects the applied participant,
whereas overt \ird{-n-} selects the patient and licenses a passive. These
contrasts show why the temporal profile shared by the three aspects cannot by
itself determine which participant is privileged (§\,\ref{sec:vm-asp}).

\subsection{The continuous (\ird{-a})}
The continuous is the unmarked member of the imperfective domain. It presents
a situation as obtaining at reference time without foregrounding change, onset,
active unfolding, or prospective projection. It is therefore the natural form
for stable states, contingent present conditions, broad durative situations, and
habitual or generic activity where what matters is that the situation holds
rather than how it is developing:

\pex
\a
\begingl
\gla Praže to mulša.//
\glb Prague.\Ind{} \textsc{loc} live-\Cont{}//
\glft \trsl{(I) live in Prague.} \lingcontext{stable durative state}//
\endgl
\a
\begingl
\gla Nužo mužna.//
\glb city.\Dir{}.\Def{} be.beautiful-\Cont{}//
\glft \trsl{The city is beautiful.} \lingcontext{property holding at the reference time}//
\endgl
\xe

Because the continuous does not foreground development or change, it is
compatible with predicates across the entire stative-dynamic spectrum, from
permanent properties to ongoing activities. Where the speaker wishes simply to
assert that a situation holds---without committing to whether it is in active
motion or stable---the continuous is the default choice. Its unmarked status
within the imperfective domain is reflected in the morphology: \ird{-a} is the
shortest and phonologically least specified of the three imperfective suffixes.

Watching television as a habit and watching it on a particular occasion
permit a contrast between the two non-projective imperfectives:

\pex\label{ex:vm-slot-habit-scene}
\a
\begingl
\gla Veň \textnormal{[[television]]}\textsubscript{\Ind{}} \textnormal{[[watch]]}\textsubscript{\Cont{}}.//
\glb woman.\Dir{} television.\Ind{} watch-\Cont{}//
\glft \trsl{A woman watches television.}
\lingcontext{a description of an evening habit; the television need not be on at speech time}//
\endgl
\a
\begingl
\gla Veň \textnormal{[[television]]}\textsubscript{\Ind{}} \textnormal{[[watch]]}\textsubscript{\Prog{}}.//
\glb woman.\Dir{} television.\Ind{} watch-\Prog{}//
\glft \trsl{A woman is watching television.}
\lingcontext{the woman is watching television at the reference point}//
\endgl
\xe

The habitual interpretation in (a) comes from the evening-routine context;
continuous morphology also permits a situation simply holding on one occasion.
The progressive in (b) selects the unfolding activity, while preserving the
same agent-direct, patient-indirect frame.

\subsection{The progressive (\ird{-ime})}
The progressive is the marked dynamic member of the imperfective domain.
Where the continuous merely asserts that a situation holds, the progressive
profiles an event specifically in its active unfolding phase---the period during
which it is dynamically occurring, not yet complete, and currently in motion.
It is accordingly most natural with activities and developing processes, where
the sense of an event in progress is pragmatically salient:

\pex
\a
\begingl
\gla Marek zemácime.//
\glb Marek.\Dir{} run-\Prog{}//
\glft \trsl{Marek is running.} \lingcontext{dynamic activity in progress}//
\endgl
\a
\begingl
\gla Marek travě pjaštíme.//
\glb Marek.\Dir{} bread.\Ind{} eat-\Prog{}//
\glft \trsl{Marek is eating bread.} \lingcontext{ongoing activity; patient in indirect}//
\endgl
\xe

The progressive of a simple state is licensed only under a marked
stage-level or expressive interpretation, such as looking particularly beautiful
on a particular occasion. It does not itself introduce becoming. Ordinary
change of state requires the overt inchoative stem in \ird{-el-}/\ird{-él-}.
The following pair contrasts a holding state with a developing transition:

\pex
\a
\begingl
\gla Nužo mužna.//
\glb city.\Dir{}.\Def{} be.beautiful-\Cont{}//
\glft \trsl{The city is beautiful.} \lingcontext{continuous: state holds at $R$}//
\endgl
\a
\begingl
\gla Nužo mužnélime.//
\glb city.\Dir{}.\Def{} be.beautiful-\textsc{inch}-\Prog{}//
\glft \trsl{The city is becoming beautiful.} \lingcontext{progressive: dynamic unfolding at $R$}//
\endgl
\xe

The continuous describes a holding situation; the progressive profiles an
unfolding eventuality whose internal structure must be independently licensed.
Thus \ird{Nužo mužna} and \ird{Nužo mužnélime} differ both in their
lower predicate and in viewpoint. Temporary bodily conditions remain states:
\ird{ptakna} means \trsl{is tired}, while an overt inchoative supplies
\trsl{become tired}. Neither bodily temporariness nor progressive aspect
alone introduces a transition (§\,\ref{sec:stative-stability}).


\subsection{The contemplative (\ird{-áž})}
Default \ird{-áž} represents /aːʒ/ in ordinary short-root
Classes~A and~B. Class~C and locally quantity-bearing valency stems
select short \ird{-až}, as in \ird{knářaž} and \ird{pjáštraž}.
Within each selected allomorph its quantity is retained before further
suffixes, while its final consonant has [ʃ] in an isolated
word by regular final devoicing. Before a following vowel, /ʒ/ retains
voicing: linked \ird{mužnáže}, evidential \ird{pjaštážal}, and
participant nominal \ird{pjaštážou} all contain [ʒ]. The linker
preserves this postalveolar consonant rather than selecting a different
stem grade.

The contemplative is the projective member of the imperfective domain. It
presents the eventuality as not yet realised at reference time but projected
from that point: the event is construed as intended, expected, impending,
planned, or otherwise prospectively oriented. The contemplative therefore does
not encode future tense---like all Iridian Asp.Align categories it is
indifferent to the relation between reference time and speech time---but rather
a prospective viewpoint in which reference time is interior to the projection
and the event lies ahead of it:

\pex
\a
\begingl
\gla Praže de stojáž.//
\glb Prah-ě de stoj-áž.//
\glc Prague.\Ind{} to go-\Ctpl{}//
\glft \trsl{(I) am about to go to Prague. / (I) intend to go to Prague.}//
\endgl
\a
\begingl
\gla Nužo mužnáž.//
\glb city.\Dir{}.\Def{} be.beautiful-\Ctpl{}//
\glft \trsl{The city will be beautiful. / The city is expected to be beautiful.}//
\endgl
\xe

A friend's leaving can likewise be projected before it begins or viewed
once it is unfolding:

\pex\label{ex:vm-slot-projection-process}
\a
\begingl
\gla Radyň varžáž.//
\glb friend.\Dir{} leave-\Ctpl{}//
\glft \trsl{A friend is about to leave.}
\lingcontext{leaving is impending but has not begun at $R$}//
\endgl
\a
\begingl
\gla Radyň váržime.//
\glb friend.\Dir{} leave-\Prog{}//
\glft \trsl{A friend is leaving.}
\lingcontext{the friend is walking out; leaving is construed as an extended process}//
\endgl
\xe

The projected leaving in (a) may subsequently be cancelled. The progressive
in (b) instead locates the reference point within the friend's movement out
of the room; neither example asserts arrival at a destination.

The projective semantics account for the range of English translations the
contemplative receives---\textit{about to}, \textit{intending to},
\textit{planning to}, or a simple future---depending on predicate type and
discourse context. With achievement predicates, which denote instantaneous
transitions, the contemplative is especially natural: instantaneous transitions
lend themselves readily to imminence and anticipation construals, where the
event is imminent or impending rather than in active progress.

The contemplative belongs to the imperfective domain because reference time
precedes the projected eventuality, whose completion is unasserted. An ordinary
active contemplative selects its agent, or its sole participant with a stative
or intransitive base; this NOM-domain orientation remains subject to the same
valency and Voice alternations as the continuous and progressive.

The three imperfective categories interact systematically with predicate class,
showing that they are distinct viewpoint operators rather than translation
variants of a single value. With {states}, the continuous gives ordinary
present-state predication; the progressive requires a marked stage-level or expressive construal;
a transition requires the overt inchoative; the contemplative gives an anticipated or expected-state reading
(\trsl{will be \ldots}, \trsl{is expected to be \ldots}). With {activities}, the
continuous gives broad ongoing or habitual interpretation; the progressive gives
a vivid, scene-internal reading in which the activity is actively in course; the
contemplative gives an intended or pending-activity reading. With
{accomplishments}, the continuous selects the internal phase without projecting
completion; the progressive highlights active advancement toward culmination;
the contemplative presents the event as pending or aimed at, with the endpoint
anticipated but not asserted. With {achievements}, which denote instantaneous
transitions, the continuous and progressive typically require coercion into
preparatory, iterative, incipient, or extended interpretations, while the
contemplative is often especially natural because instantaneous transitions
readily support imminence readings.

Imperfective clauses permit descriptions of an agent's ongoing activity,
condition, or intentions, although their use does not require an agent topic.
The three categories support distinct discourse functions. The continuous is the natural form for description and
framing---durable predications, stable conditions, and ongoing situations that
provide the setting or context for dynamic foreground events. It is the unmarked
imperfective of persistence and background. The progressive is suited to
scene-internal narration, vivid observation, and dynamic immediacy---contexts
where an event is presented as actively unfolding before the narrator's or
speaker's attention. It is characteristically a foregrounded imperfective: not
the background frame but the observed process. The contemplative serves
planning, expectation, prefiguration, and other forward-looking discourse
functions---expressing what an agent intends, what is expected to follow, or
what is imminent in the developing narrative.

Continuous and progressive predicates supply ongoing situations for phasal
comparison: \ird{so} presents their persistence, while \ird{vy} can mark
attainment of the ongoing state. Perfective predicates also permit phasal
modification under the comparison domains specified in
§\,\ref{sec:class1-particles}; viewpoint aspect supplies the predication
to which Phas applies, rather than licensing it solely through
non-boundedness.

With contemplative predicates, \ird{so} presents a projected or intended
situation as continuing, whereas \ird{vy} can mark that projection as
settled or imminent in the relevant plan. Neither combination asserts actual
completion. The \trsl{not yet} interpretation requires a negative polarised
input to \ird{so}; prospective viewpoint alone does not supply that negation
(§\,\ref{sec:class1-particles}).

\section{Mood}\label{sec:mood}

The mood system comprises the indicative, the imperative, the prohibitive, the
hortative, and the irrealis. Mood is expressed by a suffix occupying the
position immediately outside the aspect-alignment morpheme; in the indicative
this suffix is null. Where a non-indicative mood suffix is present, the
Asp.Align head remains in the clause, while Mood licenses a two-way
imperfective/perfective-domain contrast in place of the full five-way
indicative inventory. Post-syntactic impoverishment regulates the realisation
of this reduced configuration, as described in
§\,\ref{sec:vm-mood-asp}. Every finite verb, indicative or not, therefore
carries an Asp.Align exponent. The imperative, prohibitive, and hortative
together form the directive domain, a unified subsystem whose case properties
follow from directive orientation within the reduced Asp.Align domain;
it is treated as a group beginning in §\,\ref{sec:mood-directive}.

The indicative is the unmarked mood. It alone hosts the full five-way aspect
system and the associated split-ergative alignment described in
Section~\ref{sec:aspect}. Non-indicative moods carry a reduced two-way
aspect-domain contrast in Asp.Align (§\,\ref{sec:vm-mood-asp}), which collapses
the five-way viewpoint inventory; directive Mood additionally fixes the
orientation described in §\,\ref{sec:mood-directive-case}. Overt antipassive morphology is licensed only
with indicative ERG-aspect features, including when that configuration is
embedded under participant nominalisation.

\subsection{Asp.Align under non-indicative Mood: impoverishment and reduced opposition}
\label{sec:vm-mood-asp}

Every finite Iridian verb contains Asp.Align, whose exponent can be zero.
Non-indicative Mood selects a reduced aspectual configuration in syntax,
retaining an IPFV/PFV opposition together with the construction's orientation
requirements.

In the indicative, Asp.Align realises the full five-way viewpoint inventory
(progressive, continuous, contemplative, perfective, retrospective) with the
associated split-ergative alignment. The non-indicative moods---imperative,
prohibitive, hortative, and irrealis---select the coarser
\term{reduced Asp.Align} opposition shown in
Table~\ref{tab:reduced-asp}.

\begin{table}
  \footnotesize
  \caption{Reduced Asp.Align exponents in non-indicative clauses.}
  \ifparadigmappendix\else\label{tab:reduced-asp}\fi
  \begin{tabularx}{0.9\textwidth}{l l L}
    \toprule
    \textsc{domain} & \textsc{exponent} & \textsc{semantic range} \\
    \midrule
    Imperfective (IPFV) & $\varnothing$ & unbounded, ongoing, habitual, generic;
                                        processual commands and optatives \\
    Perfective (PFV)    & \ird{-ov-}/\ird{-óv-} & bounded, single whole event, completive;
                                        ingressively directive \\
    \bottomrule
  \end{tabularx}
\end{table}


Reduced Asp.Align retains the two aspectual domains, while the clause's
orientation also depends on Mood. In the irrealis, ordinary IPFV clauses retain
agent-direct and patient-indirect case, and PFV clauses retain patient-direct
and agent-transitive case. Directive Mood imposes theme orientation in both
domains (§\,\ref{sec:mood-directive-case}). The reduced temporal opposition
therefore does not require every non-indicative construction to retain the
ordinary indicative case frame.

The surface template for a non-indicative finite verb is therefore:

\begin{center}
\Stem{} + (\Appl{}) + (\Caus{}) + (\textsc{orientation}) + \textsc{reduced Asp.Align} + \textsc{mood}
\end{center}

\noindent where the reduced Asp.Align is silent ($\varnothing$) in the IPFV
domain and has \ird{-ov-} in the PFV domain, with locally strong
\ird{-óv-} selected before affirmative irrealis. Directive Mood retains
\ird{-ov-}. Representative paradigms
for the stem \ird{pjašt-} `eat':

\begin{center}
\footnotesize
\begin{tabular}{l l l l}
\toprule
\textsc{domain} & \textsc{imperative} & \textsc{prohibitive} & \textsc{hortative} \\
\midrule
IPFV ($\varnothing$) & \ird{pjaštím}   & \ird{pjašté}    & \ird{pjaštká}   \\
PFV (\ird{-ov-})   & \ird{pjaštovím} & \ird{pjaštové}  & \ird{pjaštovká} \\
\bottomrule
\end{tabular}
\end{center}

The PFV exponent \ird{-ov-} preserves the thematic consonant of the lexical
root. When irrealis \ird{-ěv} follows, its mixed-grade instruction targets
the final /v/ of the aspect exponent, whose palatal realisation is [ʋʲ],
while the intervening exponent prevents mutation of the root-final cluster.
The resulting form is \ird{pjaštóvěv}, rather than \ird{*pjáščóvěv};
the absence of root mutation is compatible with palatality at the suffixal
boundary (§\,\ref{sec:suffix-classes}).

The IPFV imperative encodes an ongoing,
habitual, or processual directive (the event is construed as
unbounded), while the PFV imperative encodes a bounded, completive
directive (the event is construed as a single whole). In the irrealis,
the IPFV domain covers generic non-actual scenarios and standing
hypotheticals; the PFV domain covers single bounded counterfactuals.

The reduced opposition distinguishes a standing instruction from a single
occasion without changing the directive case frame:

\pex\label{ex:vm-slot-directive-domains}
\a
\begingl
\gla \textnormal{[[dishes]]}\textsubscript{\Dir{}} \textnormal{[[wash]]}\textsubscript{\textsc{ipfv.imp}}.//
\glb dishes.\Dir{} wash-\textsc{ipfv}-\Imp{}//
\glft \trsl{Wash dishes regularly.}
\lingcontext{a household rule after meals; recurrence is supplied by the context}//
\endgl
\a
\begingl
\gla \textnormal{[[dishes]]}\textsubscript{\Dir{}} \textnormal{[[wash]]}\textsubscript{\textsc{pfv.imp}}.//
\glb dishes.\Dir{} wash-\textsc{pfv}-\Imp{}//
\glft \trsl{Wash the dishes.}
\lingcontext{the dishes from one meal are to be washed before going out}//
\endgl
\xe

Both instructions select the dishes as direct; the understood addressee
retains the transitive case required by a directive with a distinct theme.
The recurrence in the first translation comes from the household context:
IPFV does not itself encode a frequency, and PFV presents the commanded
operation as bounded rather than adding a separate degree modifier.

In the Distributed Morphology implementation, Vocabulary Insertion maps the
coarse IPFV/PFV distinction to $\varnothing$ or \ird{-ov-}. A PF
impoverishment rule can exclude the finer indicative exponents in this
environment, although it does not supply the syntactic reduction, determine
case, or remove semantic distinctions after interpretation. The
exponence restriction is:

\medskip
\noindent\begin{tabular}{@{}l@{\quad}l}
$[\textsc{Mood}: \text{non-indicative}]$ & $\Rightarrow$ realise only the coarse
aspectual contrast at Asp.Align; \\
& \hspace*{1.5em}retain only $[\pm\textsc{pfv-domain}]$.
\end{tabular}
\medskip

Mood is structurally higher than Asp, with intervening Neg and licensed
phasal or comment layers where present
(Table~\ref{tab:vm-full-structure}). Its syntactic selection establishes the
reduced configuration before case valuation, while its exponent follows the
zero or overt Asp.Align exponent. Morphological neutralisation therefore
presupposes the licensed configuration rather than deriving its case frame.

Directive Mood imposes an orientation requirement on the reduced aspectual
domain: a retained theme takes direct case, regardless of the IPFV/PFV value.
Where no theme is retained, selection falls on an applied participant, then a
causee, or the sole argument of an underived intransitive. The agent or causer
of a clause with a distinct direct participant takes transitive case, whether
expressed or contextually understood. This requirement is specific to the
directive construction; omission of a subject does not transfer its case to
another argument (§\,\ref{sec:mood-directive-case}).


\noindent Directive constructions and irrealis are treated in
Chapter~\ref{ch:clause-types}, while the polarity contrasts of Mood are
analysed in §\,\ref{sec:vm-moodp}.

\section{Inner evidential morphology}\label{sec:vm-ip}
\label{sec:vm-evid}

An evidential suffix occupies the outermost position in the finite inflectional
spine, immediately inside the complement linker \ird{-e}. This \term{inner
evidential} (\synhead{EvidP}$_\text{inner}$) is distinct from the clausal
evidential and speech-act operators in the left periphery
(§\,\ref{sec:th-particle-framework}), which are Zone~I particles occupying
\synhead{EvidP}$_\text{outer}$ and \synhead{SpeechActP} in the clause spine.

The indirect evidential marks the source of knowledge for the proposition
expressed by the finite verb. Its default exponent is \ird{-al}, with
\ird{-l} selected after a phonologically vowel-final host. The conditioning
depends on the final segment of the fully formed finite host, including any
mood exponent: consonant-final \ird{pjaštník} takes \ird{-al}, whereas
vowel-final \ird{pjašta} and \ird{pjaštíme} take \ird{-l}. A long vowel
or a final diphthongal nucleus also selects \ird{-l}. Selection precedes
coda vocalisation: the /v/ of retrospective \ird{-av} is a consonant
and therefore selects \ird{-al}, even though its free word-final
realisation is [w]. The host retains its vowel quantity; this attachment
requires neither contraction nor an inserted \ird{v}.

The evidential selects the plain grade at its attachment boundary. It adds
no root mutation, vocalic stem grade, or change of orientation and leaves
the finite host's aspect and mood exponents intact. It supplies neither
tense nor person, gender, or number agreement. Both allomorphs belong to
the host's lexical prosodic word and preserve its initial stress.

A report about a child singing illustrates the two allomorphs
without changing the child's case:

\pex\label{ex:vm-slot-evidential-hosts}
\a
\begingl
\gla Byl kánimel.//
\glb child.\Dir{} sing-\Prog{}-\textsc{ind.evid}//
\glft \trsl{A child is reportedly singing.}
\lingcontext{the report concerns an activity in progress; vowel-final \ird{kánime} selects \ird{-l}}//
\endgl
\a
\begingl
\gla Byl kánekal.//
\glb child.\Dir{} sing-\Pf{}-\textsc{ind.evid}//
\glft \trsl{A child reportedly sang.}
\lingcontext{the report concerns a bounded performance; consonant-final \ird{kánek} selects \ird{-al}}//
\endgl
\xe

The intransitive reading of \ird{kan-} retains a sole direct participant in
both aspects. The suffix contrast follows the host's final segment, while
the reported source of information leaves each aspectual interpretation intact.

The order of the outer suffixes is:

\begin{center}
\textsc{stem} + \synhead{Asp.Align} + (\textsc{mood}) +
(\textsc{ind.evid}) + \ird{-e} + \ird{po}
\end{center}

\noindent The corresponding syntactic projection is above
\synhead{MoodP} and constitutes the outermost layer of the finite inflectional
spine. The boundary with the clausal \synhead{EvidP}$_\text{outer}$ (a
left-peripheral projection) must be kept clear: the inner evidential marks the
\emph{source of knowledge} as it bears on the proposition expressed by the
verbal complex; the outer \synhead{EvidP} particles (Zone~I) operate at the
clause level. Their relation to interrogative force depends on the licensed
combination; ordinary pre-particle pronunciation of an argument supplies no
independent claim that it takes scope over the evidential. The discourse-relational and speech-act framing in
\synhead{SpeechActP} occupies the higher layer of the clause spine
(Table~\ref{tab:th-clause-spine}).

\begin{table}[htbp]
\footnotesize
\caption{Indirect evidential forms across the five indicative aspects.}
\ifparadigmappendix\else\label{tab:inner-evidential-paradigm}\fi
\begin{tabularx}{\textwidth}{L L L L}
\toprule
{\sc aspect} & {\sc finite host} & {\sc indirect evidential} & {\sc linked evidential} \\
\midrule
Perfective & \ird{pjaštník} & \ird{pjaštníkal} & \ird{pjaštníkalě} \\
Retrospective & \ird{pjaštnáv} & \ird{pjaštnával} & \ird{pjaštnávalě} \\
Continuous & \ird{pjašta} & \ird{pjaštal} & \ird{pjaštalě} \\
Progressive & \ird{pjaštíme} & \ird{pjaštímel} & \ird{pjaštímelě} \\
Contemplative & \ird{pjaštáž} & \ird{pjaštážal} & \ird{pjaštážalě} \\
\bottomrule
\end{tabularx}
\end{table}


The forms in Table~\ref{tab:inner-evidential-paradigm} retain the ordinary
transitive orientation appropriate to each aspect. The retrospective
contains long \ird{-náv}, so that its evidential ending is
\ird{-nával}. Strong antipassive \ird{pjáščav} instead preserves
short \ird{-av} in \ird{pjáščaval}. The \ird{a} of continuous \ird{pjašta} belongs to
Asp.Align: \ird{pjašta-l} yields \ird{pjaštal}, without deleting an
aspect vowel or adding a second one.

When \ird{-e} follows the evidential, its front-grade mapping targets
evidential-final /l/, which selects [ʎ], written \ird{lě} before /ɛ/. Thus
\ird{-al-e} gives \ird{-alě}, and postvocalic \ird{-l-e} gives
\ird{-lě}. The preceding finite host remains intact in both:
\ird{pjaštíme-l-e} yields \ird{pjaštímelě}. The \ird{e} before
\ird{lě} in the sequence \ird{-elě} is the progressive's own final vowel.
Syllabification follows the complete suffixal sequence, so a free
evidential's final hard /l/ vocalises to [w], whereas the linked soft
lateral remains [ʎ] (§\,\ref{sec:l-vocalisation}).

This outer boundary does not reach through the evidential to select
inner linker allomorphs. Perfective \ird{pjášček} accordingly forms
\ird{pjáščekal} and linked \ird{pjáščekalě}, preserving \ird{-ek};
its directly linked counterpart \ird{pjáščice} belongs to the distinct
construction without intervening Evid. Likewise, \ird{pjaštímelě}
retains Class~B \ird{-íme}, whereas direct linking without Evid selects
\ird{pjaštímně}. The contrast follows the local boundary of each
selecting exponent (§\,\ref{sec:boundary-mapping}).

The suffix attaches outside the independently licensed finite construction.
Table~\ref{tab:evidential-combinations} gives constructed combinations;
the causative forms preserve both the rhotic exponent and outer
orientation \ird{-an-} where required. Mood supplies the input to
evidential allomorphy, so prohibitive \ird{pjašté} selects \ird{-l}
while imperative \ird{pjaštím} selects \ird{-al}. Their long vowels
remain long. The directive rows illustrate exponence where a combination
of Evid and directive Mood is independently licensed; they do not establish
its semantic or pragmatic distribution. The choice of \ird{-al} or
\ird{-l} alone does not settle whether such a combination is available.

\begin{table}[htbp]
\footnotesize
\caption{Constructed combinations of the evidential with finite morphology.}
\label{tab:evidential-combinations}
\begin{tabularx}{\textwidth}{L L L L}
\toprule
{\sc construction} & {\sc finite host} & {\sc indirect evidential} & {\sc linked evidential} \\
\midrule
Antipassive perfective & \ird{pjášček} & \ird{pjáščekal} & \ird{pjáščekalě} \\
Antipassive retrospective & \ird{pjáščav} & \ird{pjáščaval} & \ird{pjáščavalě} \\
Passive progressive & \ird{pjaštníme} & \ird{pjaštnímel} & \ird{pjaštnímelě} \\
Causative progressive & \ird{pjáštrime} & \ird{pjáštrimel} & \ird{pjáštrimelě} \\
Causative perfective & \ird{pjáštranik} & \ird{pjáštranikal} & \ird{pjáštranikalě} \\
Irrealis & \ird{pjáščev} & \ird{pjáščeval} & \ird{pjáščevalě} \\
Negative irrealis & \ird{pjaštir} & \ird{pjaštiral} & \ird{pjaštiralě} \\
Imperative & \ird{pjaštím} & \ird{pjaštímal} & \ird{pjaštímalě} \\
Prohibitive & \ird{pjašté} & \ird{pjaštél} & \ird{pjaštélě} \\
\bottomrule
\end{tabularx}
\end{table}

Indicative negation likewise preserves its ordinary form outside the lexical
stem: \ird{zapjaštník} gives \ird{zapjaštníkal} and linked
\ird{zapjaštníkalě}. Evidential suffixation neither replaces negative
Mood \ird{-ir} with \ird{za-} nor changes the argument frame licensed
by the inner verb.
